\section{Working with PostgreSQL}
%% ═════════════════════════════════════════════════════════════════════════════

\begin{frame}[fragile]{\textcolor{brown}{Setting Up a mock Database}}
\begin{block}{Create and connect to the database}
\begin{lstlisting}[style=sql]
CREATE DATABASE my_database;
\c my_database
\end{lstlisting}

\end{block}

\begin{block}{Create the \texttt{students} table with a default value on \texttt{enrolled}}
\lstinputlisting[style=sql]{component/src/create_std.sql}

\end{block}

  % \begin{block}{\faLightbulb\; Note on \texttt{DEFAULT TRUE}}
  %   When \texttt{enrolled} is omitted in an \texttt{INSERT} statement, PostgreSQL
  %   automatically sets it to \texttt{TRUE}. You will see this in the next slide.
  % \end{block}
\end{frame}

% ─────────────────────────────────────────────────────────────────────────────
\begin{frame}[fragile]{\textcolor{brown}{Inserting Records (DML — INSERT)}}
\begin{block}{Insert records one at a time}
\lstinputlisting[style=sql]{component/src/insert_one.sql}

\end{block}

\begin{block}{Insert multiple records in a single query}
\lstinputlisting[style=sql]{component/src/insert_mult.sql}
\end{block}

  % \begin{alertblock}{\faBolt\; Tip}
  %   The last five records omit \texttt{enrolled} — PostgreSQL silently fills it with
  %   the \texttt{DEFAULT} value of \texttt{TRUE}.
  % \end{alertblock}
\end{frame}

% ─────────────────────────────────────────────────────────────────────────────
\begin{frame}[fragile]{\textcolor{brown}{Reading Data (DML — SELECT)}}

\begin{block}{Reading records from a table}
\lstinputlisting[style=sql]{component/src/select.sql}
\end{block}
\end{frame}

% ─────────────────────────────────────────────────────────────────────────────
\begin{frame}[fragile]{\textcolor{brown}{Updating \& Deleting Records}}
  \begin{block}{\textbf{UPDATE} — change an existing value}
\lstinputlisting[style=sql]{component/src/update.sql}
\end{block}

  \vspace{0.3em}
\begin{block}{\textbf{DELETE} — remove a specific row}
\lstinputlisting[style=sql]{component/src/delete.sql}
\end{block}

  \vspace{0.3em}
  % \begin{alertblock}{\faExclamationTriangle\; Critical Rule}
  \emph{\textbf{NB: }}\textbf{Always include a \texttt{WHERE} clause} with both \texttt{UPDATE} and \texttt{DELETE}.
    Without it, \\[4pt]
  \textbf{every row} in the table will be modified or deleted.
  % \end{alertblock}
\end{frame}